\documentclass[8pt,letterpaper]{extarticle}
\usepackage[top=0.3in,bottom=0.3in,left=0.3in,right=0.3in]{geometry}
\usepackage{multicol}
\usepackage{amsmath,amssymb}
\usepackage{booktabs}
\usepackage{array}
\usepackage{enumitem}
\usepackage{xcolor}
\usepackage{titlesec}
\usepackage{parskip}


\setlength{\parskip}{1pt}
\setlength{\parindent}{0pt}
\setlength{\columnsep}{8pt}
\setlength{\columnseprule}{0.3pt}

\titleformat{\section}[block]{\normalsize\bfseries\centering\color{blue!70!black}}{}{0em}{}[\vspace{-2pt}\rule{\linewidth}{0.5pt}\vspace{-4pt}]
\titleformat{\subsection}[runin]{\small\bfseries\color{red!60!black}}{}{0em}{}[:\quad]
\titlespacing{\section}{0pt}{3pt}{1pt}
\titlespacing{\subsection}{0pt}{2pt}{0pt}

\newcommand{\prob}[1]{\smallskip\noindent{\footnotesize\bfseries\color{black!80} #1\quad}}
\newcommand{\soln}{\noindent{\footnotesize\itshape\color{blue!50!black}Sol: }}
\newcommand{\ans}[1]{{\footnotesize\bfseries\color{green!40!black}$\Rightarrow$ #1}}

\begin{document}
\begin{multicols*}{3}

%=============================================================
\section*{\small ENGR 232 — All Homework Problems \& Solutions}
%=============================================================

% ==================== HOMEWORK 1 ====================
\section*{HW1 — Ch.2: Atomic Structure \& Bonding}

\prob{2.2} Cr has 4 isotopes: 4.34\% of $^{50}$Cr (49.9460 amu), 83.79\% of $^{52}$Cr (51.9405 amu), 9.50\% of $^{53}$Cr (52.9407 amu), 2.37\% of $^{54}$Cr (53.9389 amu). Confirm avg atomic weight = 51.9963 amu.

\soln $\bar{A} = \sum f_i A_i$:
\[
  \bar{A} = (0.0434)(49.9460)+(0.8379)(51.9405)
\]
\[
  +(0.0950)(52.9407)+(0.0237)(53.9389)
\]
\ans{$= 51.9963$~amu \checkmark}

\prob{2.5} (a) Grams in 1 amu? (b) Atoms in 1 lb-mol?

\soln (a) $\dfrac{1\text{ g/mol}}{6.022\times10^{23}\text{ atoms/mol}} = $\ans{$1.66\times10^{-24}$~g/amu}

(b) $453.6\text{ g/lb}_m \times 6.022\times10^{23}\text{ atoms/mol} = $\ans{$2.73\times10^{26}$~atoms/lb-mol}

\prob{2.15} Force of attraction between K$^+$ and O$^{2-}$ ions separated by 1.5~nm.

\soln $F_A = \dfrac{A}{r^2}$, where $A = \dfrac{Z_1 Z_2 e^2}{4\pi\varepsilon_0}$, $Z_1=1$, $Z_2=2$, $r=1.5$~nm
\ans{$F_A = 2.05\times10^{-10}$~N}

\prob{2.22} (a) Main differences between ionic, covalent, metallic bonding? (b) Pauli exclusion principle?

\soln (a) \textit{Ionic}: electrostatic attraction between oppositely charged ions. \textit{Covalent}: electron sharing so each atom achieves stable config. \textit{Metallic}: ion cores shielded/glued by sea of valence electrons.
(b) Each electron state holds $\leq2$ electrons with opposite spins.

\prob{2.25} \% ionic character of TiO$_2$, ZnTe, CsCl, InSb, MgCl$_2$.

\soln $\%IC = \left[1-e^{-0.25(X_A-X_B)^2}\right]\times100$

TiO$_2$: $X_{Ti}=1.5$, $X_O=3.5$ \ans{$\approx63.2\%$}\\
ZnTe: $X_{Zn}=1.6$, $X_{Te}=2.1$ \ans{$\approx6.1\%$}\\
CsCl: $X_{Cs}=0.7$, $X_{Cl}=3.0$ \ans{$\approx74.3\%$}\\
InSb: $X_{In}=1.7$, $X_{Sb}=1.9$ \ans{$\approx1.0\%$}\\
MgCl$_2$: $X_{Mg}=1.2$, $X_{Cl}=3.0$ \ans{$\approx55.5\%$}

% ==================== HOMEWORK 2 ====================
\section*{HW2 — Ch.3: Crystal Structures}

\prob{3.3} Show for BCC that $a = \dfrac{4R}{\sqrt{3}}$.

\soln Using body diagonal of BCC unit cell: the diagonal $= 4R$ and $= a\sqrt{3}$. Therefore:
\ans{$a = \dfrac{4R}{\sqrt{3}}$} \checkmark

\prob{3.7} Density of iron (BCC, $R=0.124$~nm, $A=55.85$~g/mol).

\soln $\rho = \dfrac{nA}{V_C N_A}$, BCC: $n=2$, $V_C = a^3 = \left(\dfrac{4R}{\sqrt{3}}\right)^3$
\[a = \dfrac{4(0.124)}{\sqrt{3}} = 0.2866\text{ nm}\]
\ans{$\rho = 7.90$~g/cm$^3$} (exp: 7.87~g/cm$^3$)

\prob{3.24} Point indices for all atoms in BCC unit cell.

\soln Corner atoms at: 000, 100, 110, 010, 001, 101, 111, 011. Body-center atom:
\ans{$\frac{1}{2}\frac{1}{2}\frac{1}{2}$}

\prob{3.43} Indices for two crystallographic planes (one parallel to $x,z$, intercepts $y$ at $b/2$; one intercepts at $a/2$, $-b/2$, $c$).

\soln Plane 1: $A=\infty a$, $B=b/2$, $C=\infty c$ $\Rightarrow h=0, k=2, l=0$ \ans{(020)}

Plane 2: $A=a/2$, $B=-b/2$, $C=c$ $\Rightarrow h=2, k=-2, l=1$ \ans{$(2\bar{2}1)$}

\prob{3.63} Interplanar spacing $d_{111}$ for Mo (BCC, $R=0.1363$~nm).

\soln $a = \dfrac{4R}{\sqrt{3}} = \dfrac{4(0.1363)}{\sqrt{3}} = 0.3148$~nm

$d_{hkl} = \dfrac{a}{\sqrt{h^2+k^2+l^2}}$, so $d_{111} = \dfrac{0.3148}{\sqrt{3}}$ \ans{$= 0.1817$~nm}

% ==================== HOMEWORK 3 ====================
\section*{HW3 — Ch.4: Imperfections in Solids}

\prob{4.2} Hypothetical metal: $N_v=2.8\times10^{24}$~m$^{-3}$ at 750°C, $\rho=5.60$~g/cm$^3$, $A=65.6$~g/mol. Find vacancy fraction.

\soln $N = \dfrac{N_A \rho}{A} = \dfrac{(6.022\times10^{23})(5.60\times10^6)}{65.6}$
$= 5.14\times10^{28}$~m$^{-3}$

\ans{$N_v/N = \dfrac{2.8\times10^{24}}{5.14\times10^{28}} = 5.44\times10^{-5}$}

\prob{4.4} Vacancies/m$^3$ in Fe at 850°C ($Q_v=1.08$~eV, $\rho=7.65$~g/cm$^3$, $A=55.85$~g/mol).

\soln $N_v = N\exp\left(-\dfrac{Q_v}{kT}\right)$, $T=1123$~K, $N=8.24\times10^{28}$~m$^{-3}$
\ans{$N_v = 1.18\times10^{24}$~m$^{-3}$}

\prob{4.14} Composition (wt\%) of 6 at\% Pb -- 94 at\% Sn alloy.

\soln $C_{Pb}' = \dfrac{C_{Pb}A_{Pb}}{C_{Pb}A_{Pb}+C_{Sn}A_{Sn}}\times100$

$A_{Pb}=207.2$, $A_{Sn}=118.71$ \ans{$C_{Pb}=10.0$~wt\%, $C_{Sn}=90.0$~wt\%}

\prob{4.16} Composition (at\%) of alloy with 98~g Sn + 65~g Pb.

\soln $n'_{Sn}=98/118.71=0.825$~mol, $n'_{Pb}=65/207.2=0.314$~mol
\ans{$C_{Sn}=72.4$~at\%, $C_{Pb}=27.6$~at\%}

\prob{4.26} Au atoms/cm$^3$ in 10~wt\% Au -- 90~wt\% Ag alloy ($\rho_{Au}=19.32$, $\rho_{Ag}=10.49$~g/cm$^3$).

\soln Using Equation 4.21:
$C_{Au}'' = \dfrac{C_{Au}N_A}{\dfrac{C_{Au}A_{Au}}{\rho_{Au}}+\dfrac{C_{Ag}A_{Ag}}{\rho_{Ag}}}$
\ans{$= 3.36\times10^{21}$~atoms/cm$^3$}

% ==================== HOMEWORK 4 ====================
\section*{HW4 — Ch.5: Diffusion}

\prob{5.8} kg of H$_2$ passing/hr through Pd sheet: thickness 5~mm, area 0.20~m$^2$, 500°C, $D=1.0\times10^{-8}$~m$^2$/s, $C_A=2.4$~kg/m$^3$, $C_B=0.6$~kg/m$^3$.

\soln $J = -D\dfrac{\Delta C}{\Delta x}$, $M = JAt$

$J = (1.0\times10^{-8})\dfrac{2.4-0.6}{5\times10^{-3}} = 3.6\times10^{-6}$~kg/m$^2$/s

$M = J\times A\times 3600\text{ s}$ \ans{$= 2.6\times10^{-3}$~kg/h}

\prob{5.15} N concentration at 1~mm depth after 10~h in Fe at 700°C ($C_s=0.1$~wt\%, $D=2.5\times10^{-11}$~m$^2$/s).

\soln $C_x = C_s\left[1-\text{erf}\!\left(\dfrac{x}{2\sqrt{Dt}}\right)\right]$

$\dfrac{x}{2\sqrt{Dt}} = \dfrac{10^{-3}}{2\sqrt{(2.5\times10^{-11})(36000)}} = 0.527$

erf(0.527) $\approx 0.5436$ \ans{$C_x = 0.046$~wt\% N}

\prob{5.19} $D$ for Zn in Cu at 650°C ($D_0=2.4\times10^{-5}$~m$^2$/s, $Q_d=189{,}000$~J/mol).

\soln $D = D_0\exp\!\left(-\dfrac{Q_d}{RT}\right)$, $T=923$~K
\ans{$D = 4.8\times10^{-16}$~m$^2$/s}

\prob{5.24} Ag in Cu: $D_{650}=5.5\times10^{-16}$, $D_{900}=1.3\times10^{-13}$~m$^2$/s. (a) Find $D_0$, $Q_d$. (b) $D$ at 875°C.

\soln (a) Using two equations $\ln D = \ln D_0 - Q_d/RT$:
\ans{$Q_d = 196{,}700$~J/mol}, \ans{$D_0=7.5\times10^{-5}$~m$^2$/s}

(b) $T=1148$~K: \ans{$D = 8.3\times10^{-14}$~m$^2$/s}

% ==================== HOMEWORK 5 ====================
\section*{HW5 — Ch.6: Mechanical Properties of Metals}

\prob{6.3} Elastic strain in Al specimen (10$\times$12.7~mm cross section) under 35,500~N.

\soln $\sigma = F/A_0 = 35500/(127\times10^{-6}) = 279.5$~MPa

$\varepsilon = \sigma/E = 279.5\times10^6/(69\times10^9)$ \ans{$= 4.05\times10^{-3}$}

\prob{6.6} Elongation of Ti wire ($d=3.0$~mm, $L_0=2.5\times10^4$~mm) under 500~N.

\soln $\Delta l = \dfrac{FL_0}{A_0 E}$, $A_0 = \pi(1.5\times10^{-3})^2 = 7.07\times10^{-6}$~m$^2$, $E_{Ti}=107$~GPa

\ans{$\Delta l = 16.6$~mm}

\prob{6.10} Elongation of steel alloy specimen ($d=10$~mm, $L_0=75$~mm) under 20,000~N.

\soln $\sigma = F/A_0 = 20000/(\pi(5\times10^{-3})^2) = 255$~MPa

From stress-strain curve, $\varepsilon \approx 0.0012$ (elastic region)

$\Delta l = \varepsilon L_0 = 0.0012\times75$ \ans{$= 0.090$~mm}

\prob{6.28} Steel bar ($L_0=300$~mm, sq. cross-section 4.5~mm side). (a) Load for $\Delta l=0.45$~mm. (b) Deformation after release.

\soln (a) $\varepsilon=\Delta l/L_0=0.45/300=0.0015$. From Fig.~6.22: $\sigma\approx300$~MPa (elastic).

$F = \sigma A_0 = 300\times10^6\times(4.5\times10^{-3})^2$ \ans{$= 6075$~N}

(b) \ans{No permanent deformation} (elastic only)

% ==================== HOMEWORK 6 ====================
\section*{HW6 — Ch.7 \& 8: Dislocations, Fracture, Fatigue}

\prob{7.14} Critical resolved shear stress for Ag: tensile along [001], slip on (111) in $[\bar{1}01]$, initiated at $\sigma=1.1$~MPa.

\soln $\lambda=$ angle([001],$[\bar{1}01]$): $\cos\lambda = \dfrac{(0)(-1)+(0)(0)+(1)(1)}{\sqrt{1}\sqrt{2}} = \dfrac{1}{\sqrt{2}}$

$\phi=$ angle([001],[111]): $\cos\phi = \dfrac{1}{\sqrt{3}}$

$\tau_{crss} = \sigma\cos\phi\cos\lambda = 1.1\cdot\dfrac{1}{\sqrt{3}}\cdot\dfrac{1}{\sqrt{2}}$ \ans{$= 0.45$~MPa}

\prob{7.20} Why are small-angle grain boundaries less effective at blocking slip than high-angle boundaries?

\soln \ans{Less crystallographic misalignment at small-angle boundaries $\Rightarrow$ less change in slip direction $\Rightarrow$ less interference.}

\prob{7.30} Pre-cold-worked radius of Cu wire ($r_d=10$~mm, ductility 25\%EL).

\soln From Fig.~7.19c: 25\%EL $\Rightarrow \approx11$\%CW.

$\%CW = \left(\dfrac{r_0^2-r_d^2}{r_0^2}\right)\times100 \Rightarrow 0.11 = 1-(10/r_0)^2$

\ans{$r_0 = 10/\sqrt{0.89} = 10.6$~mm}

\prob{8.4} Max allowable surface crack length in polystyrene ($\sigma=1.25$~MPa, $\gamma_s=0.50$~J/m$^2$, $E=3.0$~GPa).

\soln $a = \dfrac{2E\gamma_s}{\pi\sigma^2} = \dfrac{2(3.0\times10^9)(0.50)}{\pi(1.25\times10^6)^2}$ \ans{$= 6.1\times10^{-4}$~m $= 0.61$~mm}

\prob{8.22} Brass fatigue data. (a) S-N plot. (b) Fatigue strength at $5\times10^5$ cycles. (c) Fatigue life at 200~MPa.

\begin{center}
\footnotesize
\begin{tabular}{@{}cc@{}}
\toprule
$\sigma_a$ (MPa) & Cycles\\
\midrule
310 & $2\times10^5$\\
223 & $1\times10^6$\\
168 & $1\times10^7$\\
143 & $1\times10^8$\\
127 & $1\times10^9$\\
\bottomrule
\end{tabular}
\end{center}
(b) \ans{$\approx250$~MPa at $5\times10^5$ cycles}\\
(c) \ans{$\approx2\times10^6$ cycles at 200~MPa}

% ==================== HOMEWORK 7 ====================
\section*{HW7 — Ch.9: Phase Diagrams I}

\prob{9.11} Is it possible to have a Cu-Ni alloy in equilibrium with liquid (20~wt\%~Ni) and $\alpha$ (37~wt\%~Ni)?

\soln \ans{No.} A single tie line in the $\alpha+L$ region cannot intersect the phase boundaries at both 20~wt\%~Ni (liquidus $\approx$1200°C) and 37~wt\%~Ni (solidus $\approx$1230°C).

\prob{9.14} 50~wt\%~Pb–50~wt\%~Mg alloy cooled from 700°C to 400°C. (a) $T$ of first solid? (b) Composition of first solid? (c) $T$ of complete solidification? (d) Composition of last liquid?

\soln (a) \ans{$\approx560$°C} (intersection of vertical line with liquidus)\\
(b) \ans{21~wt\%~Pb--79~wt\%~Mg} (tie line at 560°C, solidus intercept)\\
(c) \ans{$\approx465$°C} (eutectic isotherm)\\
(d) \ans{$\approx67$~wt\%~Pb--33~wt\%~Mg} (eutectic composition)

% ==================== HOMEWORK 8 ====================
\section*{HW8 — Ch.9: Phase Diagrams II}

\prob{9.19} 90~wt\%~Ag–10~wt\%~Cu in $\beta+L$ region, liquid at 85~wt\%~Ag. (a)~$T$? (b) $\beta$ composition? (c) Mass fractions?

\soln (a) Tie line at 85~wt\%~Ag on liquidus: \ans{$\approx850$°C}

(b) Solidus intersection: \ans{$\approx95$~wt\%~Ag}

(c) Lever rule, $C_0=90$, $C_L=85$, $C_\beta=95$~wt\%~Ag:
$W_\beta = (90-85)/(95-85)$ \ans{$= 0.5$}, $W_L=0.5$

\prob{9.22} A–B alloy: $C_0=55$~wt\%~B, $W_\alpha=W_\beta=0.5$, $C_\beta=90$~wt\%~B. Find $C_\alpha$.

\soln $W_\alpha = (C_\beta-C_0)/(C_\beta-C_\alpha) = 0.5$

$(90-55)/(90-C_\alpha)=0.5 \Rightarrow 90-C_\alpha=70$ \ans{$C_\alpha=20$~wt\%~B}

\prob{9.53} Carbon concentration for total ferrite fraction = 0.94.

\soln $W_\alpha = (0.76-C_0)/(0.76-0.022)$, solve: $0.94\times0.738=0.76-C_0$

\ans{$C_0\approx0.067$~wt\%~C}

\prob{9.58} Mass fractions of proeutectoid ferrite and pearlite in 0.25~wt\%~C alloy.

\soln $W_p = \dfrac{C_0'-0.022}{0.74} = \dfrac{0.25-0.022}{0.74}$ \ans{$= 0.308$}

$W_{\alpha'} = \dfrac{0.76-C_0'}{0.74} = \dfrac{0.76-0.25}{0.74}$ \ans{$= 0.689$}

\prob{9.61} Alloy with $W_{\alpha'}=0.20$, $W_p=0.80$: find carbon concentration.

\soln From $W_p = (C_0-0.022)/0.74 = 0.80$:

$C_0 - 0.022 = 0.592$ \ans{$C_0=0.61$~wt\%~C}

% ==================== HOMEWORK 9 ====================
\section*{HW9 — Ch.10 \& 12: Phase Trans. \& Ceramics}

\prob{10.3} Critical radius for homogeneous nucleation in Cu at 849°C ($\Delta H_f=-1.77\times10^9$~J/m$^3$, $\gamma=0.200$~J/m$^2$).

\soln $r^* = -\dfrac{2\gamma}{\Delta G_v} = \dfrac{2\gamma T_m}{|\Delta H_f|(T_m-T)}$

$= \dfrac{2(0.200)(1358)}{1.77\times10^9(1358-1122)}$ \ans{$= 1.29\times10^{-9}$~m $= 1.29$~nm}

\prob{10.8} Avrami kinetics: $n=2.5$, $y=0.40$ at $t=200$~min. Find recrystallization rate.

\soln $y=1-e^{-kt^n}$, solve for $k$:

$k = -\dfrac{\ln(1-0.40)}{200^{2.5}} = \dfrac{0.5108}{1.789\times10^6} = 2.855\times10^{-7}$

$t_{0.5} = \left(\dfrac{\ln 2}{k}\right)^{1/n} = \left(\dfrac{0.693}{2.855\times10^{-7}}\right)^{0.4}$ \ans{$\approx 214$~min}

Rate $= 1/t_{0.5}$\ans{$\approx 4.7\times10^{-3}$~min$^{-1}$}

\prob{12.16} Density of FeO (rock salt structure, $r_{Fe^{2+}}=0.077$~nm, $r_{O^{2-}}=0.140$~nm).

\soln $n'=4$, $a=2(r_{Fe^{2+}}+r_{O^{2-}})=2(0.217)=0.434$~nm

$\rho = \dfrac{n'(\sum A)}{V_C N_A} = \dfrac{4(55.85+16.00)}{(4.34\times10^{-8})^3(6.022\times10^{23})}$ \ans{$= 5.84$~g/cm$^3$}

\prob{12.18} CdS: $a=0.582$~nm, $\rho=4.82$~g/cm$^3$. How many Cd$^{2+}$ and S$^{2-}$ per unit cell?

\soln $n' = \dfrac{\rho V_C N_A}{\sum A} = \dfrac{4.82(0.582\times10^{-7})^3(6.022\times10^{23})}{112.41+32.06}$

\ans{$n'\approx4$ Cd$^{2+}$ and 4 S$^{2-}$ per unit cell}

\prob{12.29} Schottky defect fraction in NaCl at 801°C ($Q_s=2.3$~eV).

\soln $\dfrac{N_s}{N} = \exp\!\left(-\dfrac{Q_s}{2kT}\right)$, $T=1074$~K, $k=8.62\times10^{-5}$~eV/K

$= \exp\!\left(-\dfrac{2.3}{2(8.62\times10^{-5})(1074)}\right)$ \ans{$= 4.03\times10^{-6}$}

% ==================== HOMEWORK 10 ====================
\section*{HW10 — Ch.14 \& 18: Polymers \& Electrical}

\prob{14.5} Polypropylene MW data: compute (a) $\bar{M}_n$, (b) $\bar{M}_w$, (c) degree of polymerization.

\begin{center}
\footnotesize
\begin{tabular}{@{}ccccc@{}}
\toprule
Range & $M_i$ & $x_i$ & $x_iM_i$ & $w_iM_i$\\
\midrule
8k--16k & 12k & 0.05 & 600 & 240\\
16k--24k & 20k & 0.16 & 3200 & 2000\\
24k--32k & 28k & 0.24 & 6720 & 5600\\
32k--40k & 36k & 0.28 & 10080 & 10800\\
40k--48k & 44k & 0.20 & 8800 & 11880\\
48k--56k & 52k & 0.07 & 3640 & 5720\\
\bottomrule
\end{tabular}
\end{center}

(a) $\bar{M}_n = \sum x_i M_i$ \ans{$= 33{,}040$~g/mol}

(b) $\bar{M}_w = \sum w_i M_i$ \ans{$= 36{,}240$~g/mol}

(c) Polypropylene repeat unit: $m = 3(12.01)+6(1.008)=42.08$~g/mol

$DP = \bar{M}_n/m = 33040/42.08$ \ans{$\approx785$}

\prob{18.3} Max length of 4-mm-dia Al wire for $R\leq2.5~\Omega$ ($\sigma_{Al}=3.8\times10^7$~$(\Omega$-m)$^{-1}$).

\soln $R=\dfrac{l}{\sigma A}$, $A=\pi(2\times10^{-3})^2=1.257\times10^{-5}$~m$^2$

$l = R\sigma A = 2.5\times3.8\times10^7\times1.257\times10^{-5}$ \ans{$= 1194$~m}

\prob{18.5} Cu wire: $d=3$~mm, $L=2$~m, $\sigma_{Cu}=6.0\times10^7$~$(\Omega$-m)$^{-1}$.

\soln $A=\pi(1.5\times10^{-3})^2=7.07\times10^{-6}$~m$^2$

(a) $R = L/(\sigma A) = 2/(6\times10^7\times7.07\times10^{-6})$ \ans{$=4.7\times10^{-3}~\Omega$}

(b) $I=V/R=0.05/4.7\times10^{-3}$ \ans{$=10.6$~A}

(c) $J=I/A=10.6/7.07\times10^{-6}$ \ans{$=1.5\times10^6$~A/m$^2$}

(d) $\mathcal{E}=V/L=0.05/2$ \ans{$=0.025$~V/m}

\prob{18.10} Drift velocity of electrons in Ge at 1000~V/m ($\mu_e=0.38$~m$^2$/V-s). Time to traverse 25~mm?

\soln (a) $v_d = \mu_e\mathcal{E} = 0.38\times1000$ \ans{$=380$~m/s}

(b) $t = l/v_d = 0.025/380$ \ans{$=6.6\times10^{-5}$~s}

\prob{18.21} Intrinsic carrier concentration of PbTe at room temp ($\sigma=500$~$(\Omega$-m)$^{-1}$, $\mu_e=0.16$, $\mu_h=0.075$~m$^2$/V-s).

\soln $\sigma = n_i e(\mu_e+\mu_h)$

$n_i = \dfrac{\sigma}{e(\mu_e+\mu_h)} = \dfrac{500}{(1.6\times10^{-19})(0.16+0.075)}$ \ans{$= 1.33\times10^{22}$~m$^{-3}$}

\vfill
\begin{center}
\footnotesize\textit{ENGR 232 — Spring 2026 — All HW Problems \& Solutions}
\end{center}

\end{multicols*}
\end{document}